% УВОД — без номер. Ориентировъчен обем: 2–3 стр.
% По А.3: „резюме на дейностите, извършени в процеса на дипломното проектиране“.
%
% КЛЮЧОВО ЗА ЦЯЛАТА РАБОТА: изходната електронна таблица НЕ е образец за
% пренаписване, а ЕТАЛОН (оракул) за коректност. Тази смяна на ролята е
% разликата между отчет „направихме уеб приложение“ и инженерен труд, в който
% има проверимо твърдение.
\section*{УВОД}
\label{sec:introduction}
\addcontentsline{toc}{section}{УВОД}

Изкуствените катерачни съоръжения се закрепват към носещата конструкция на вече
съществуваща сграда. Преди да започне проектирането на самото съоръжение, трябва да
се отговори на един въпрос, който предхожда всичко останало: издържа ли сградата
товарите, които съоръжението ще ѝ предаде. Отговорът се дава на много ранен етап,
когато за сградата се знаят само няколко величини, височина на катерачната
повърхност, разстояние между колоните и наклон на стената, и служи за преценка дали
проектът изобщо е осъществим на този адрес.

Стойностите за този отговор се съдържат в набор от таблици, изготвени от
конструктивния отдел на Walltopia и разпространявани като работна книга на Microsoft
Excel с макроси, с обем от 22{,}0 MB. Работната книга изпълнява предназначението си,
но носи ограничения, които не произтичат от съдържанието ѝ, а от носителя: изисква
конкретен настолен програмен продукт с разрешени макроси, разпространява се като
прикачен файл и затова съществува в неизвестен брой копия с неизвестни версии, не
пази следа кой, кога и с какви входни данни е получил дадено число, и не е използваема
на строителна площадка или от телефон. Отделно от това четенето ѝ изисква навик:
търсената стойност се намира на пресечната точка на три независими избора в таблица,
която не подсказва грешния избор.

\textbf{Целта} на дипломната работа е да се проектира и реализира програмна система,
която прави съдържанието на тези таблици достъпно през браузър, дава \emph{доказано
същите} числа като работната книга, запазва всяко изчисление като именуван проект на
конкретен потребител и позволява от този проект да бъде зададен въпрос към инженерния
отдел заедно с пълния изчислителен контекст.

За постигането на целта са формулирани следните \textbf{задачи}:

\begin{enumerate}[topsep=0pt,itemsep=0pt,leftmargin=*]
\item описание на проблемната област: какво описват таблиците, кои са входните и
      изходните величини, кои норми стоят зад тях, какви са недостатъците на
      съществуващото решение и какви изисквания следват от тях;
\item проектиране на архитектурата на системата и на базата данни, включително
      обосновка на избора между разгледаните архитектурни варианти, договор на
      приложния програмен интерфейс и модел на съхраняваните данни;
\item програмна реализация на двете изпълними части, услуга и уеб клиент, върху общ
      приложен програмен интерфейс, заедно с проверка на реализацията срещу еталонните
      таблици и срещу изискването интерфейсът да остава използваем при всяка ширина на
      екрана;
\item изготвяне на ръководство за потребителя, което описва работата със системата от
      гледна точка на човека, който я използва, а не на този, който я поддържа.
\end{enumerate}

\textbf{Обхватът} на работата е програмната система, тоест изчислителното ядро,
уеб интерфейсите към него и услугата, която пази потребителските данни. Извън обхвата
остават конструктивното изчисление, което е произвело самите таблици, окончателното
оразмеряване на закрепванията и всяка проверка, която подлежи на самостоятелно
конструктивно задание. Извън обхвата остава и отделното мобилно приложение, което
съществува в същото хранилище: то изостава от уеб приложението по функционалност и не
е предмет на настоящата работа. Системата дава \term{предварителни} товари, което е
записано в самия интерфейс на всяка страница с резултати и е отделено като изрично
ограничение в §\ref{subsec:limits}.

\textbf{Методологична бележка.} Работата изхожда от наблюдението, че за тази задача
съществува нещо, което рядко е налично при разработката на приложен софтуер:
\term{оракул}, тоест независим източник, за който е прието, че дава верния отговор.
Еталонната работна книга е точно такъв източник. Затова коректността на
изчислителното ядро не се твърди, а се измерва: генерира се пълното декартово
произведение на допустимите входни комбинации и всяка получена стойност се сравнява
с еталона. Същият принцип е приложен и към интерфейса, където ролята на еталон играе
не второ мнение, а механично проверимо свойство, отсъствието на хоризонтално
препълване при зададена ширина на изгледа. Проверката и получените от нея резултати
са изложени в §\ref{subsec:method}.

\textbf{Структура на изложението.} Раздел~\ref{sec:domain} описва проблемната област:
съдържанието на еталонните таблици, недостатъците на съществуващото решение,
потребителите и изведените от тях изисквания. Раздел~\ref{sec:architecture}
представя проектирането на архитектурата на системата и на базата данни, заедно с
обосновката на направените избори и договора на приложния програмен интерфейс.
Раздел~\ref{sec:implementation} излага особеностите на програмната реализация,
включително проверката на реализацията и получените резултати.
Раздел~\ref{sec:userguide} съдържа ръководството на потребителя. Заключението обобщава
предимствата и недостатъците на решението и очертава насоките за по-нататъшна работа.
